\section{Diferencias entre SMS y SLR}

Aunque ambos métodos son procesos sistemáticos y reproducibles, difieren principalmente en su alcance y profundidad, ya que el SMS se interesa en una comprensión de alto nivel como lo son el estado del arte y un estudio bibliográfico de un área más amplia, mientras que la SLR busca un estudio profundo de un tema en particular, para así responder preguntas puntuales y específicas.\\
Los resultados de una Revisión Sistemática de Literatura buscan sintetizar la evidencia de manera rigurosa para llegar a un análisis cualitativo y cuantitativo de los datos para responder a una hipótesis, en comparación con los resultados del Mapeo Sistemático que suelen ser gráficos de burbuja y mapas conceptuales.\\
Estas diferencias han sido descritas en detalle en la literatura metodológica sobre SLR y SMS en ingeniería de software. \cite{kitchenham2009slr}\\

\begin{table}[h!]
\centering
\caption{Principales diferencias entre la metodología SMS y la SLR}
\begin{tabular}{|p{2.1cm}|p{3.4cm}|p{3.4cm}|}
\hline
\textbf{Característica} & \textbf{Estudio de Mapeo Sistemático (SMS)} & \textbf{Revisión Sistemática de Literatura (SLR)} \\ \hline
\textbf{Objetivo} & Clasificar y categorizar un campo de investigación e identificar referencias disponibles. & Agregar y sintetizar evidencia para responder preguntas específicas. \\ \hline
\textbf{Alcance} & Amplio ya que busca mapear un área temática o fenómeno de interés de manera extensiva. & Específico: se centra de manera profunda en un tema o hipótesis muy puntual. \\ \hline
\textbf{Profundidad} & El análisis no es tan exhaustivo debido al volumen de estudios y el objetivo general. & Busca un entendimiento detallado y crítico de los hallazgos de los estudios primarios. \\ \hline
\textbf{Preguntas de Investigación} & Generales orientadas a tendencias, foros de publicación y distribución temática. & Puntuales y específicas como lo pueden ser factores de costo o riesgos. \\ \hline
\textbf{Evaluación de Calidad} &  Menos detallada ya que algunos criterios de calidad pueden ser menos aplicables que en una SLR. & Obligatoria y rigurosa, es esencial para asegurar que los hallazgos sean confiables y con menor sesgo. \\ \hline
\textbf{Resultados} & Mapas, tablas, gráficos de distribución y descubrimiento de brechas. & Una síntesis de evidencia que responde a la pregunta de investigación o hipótesis. \\ \hline
\textbf{Propósito estratégico} & Identificar áreas con exceso de publicaciones. & Proporcionar una base para la toma de decisiones informada o práctica basada en evidencia. \\ \hline
\end{tabular}
\end{table}